% !TeX root = main.tex

\chapter{Aplicaciones: De la Teoría a la Práctica}

\begin{quotation}
	\itshape
	``La geometría de la información no es solo matemática elegante; es la herramienta que convierte la incertidumbre en decisión.''
\end{quotation}

En este capítulo cerramos el arco del libro mostrando cómo la maquinaria geométrica desarrollada en los capítulos anteriores se aplica a problemas reales.  Cuatro áreas convergen: la \textbf{teoría de seguros} (credibilidad), el \textbf{aprendizaje automático} (optimización y generalización), la \textbf{inferencia variacional} (aproximación bayesiana) y los \textbf{métodos de gradiente natural} (optimización en espacio estadístico).

\section{Credibilidad actuarial y geometría}

\subsection{El problema de la credibilidad}

En seguros, un asegurador debe fijar primas basándose en el historial de siniestros de un asegurado.  El problema de \textbf{credibilidad} es: ¿cuánto peso dar a los datos propios del asegurado vs.\ la experiencia colectiva?

\begin{definition}[Prima de credibilidad]
	La \textbf{prima de credibilidad} es una combinación ponderada:
	\[
	P = Z \cdot \bar{X} + (1-Z) \cdot \mu
	\]
	donde $\bar{X}$ es el promedio de siniestros del asegurado, $\mu$ es la media colectiva y $Z \in [0,1]$ es el \textbf{factor de credibilidad}.
\end{definition}

\begin{theorem}[Fórmula de Bühlmann]
	En el modelo Bühlmann, el factor de credibilidad óptimo es:
	\[
	Z = \frac{n}{n + K}
	\]
	donde $K = \text{EPV}/\text{VHM}$, con EPV = varianza esperada de la pérdida y VHM = varianza de la media hipotética.
\end{theorem}

\begin{example}[Credibilidad en Poisson--Gamma]
	Para siniestros $\sim \mathrm{Poisson}(\lambda)$ con $\lambda \sim \mathrm{Gamma}(\alpha, \beta)$ (modeloBayesiano conjugado), la credibilidad es:
	\[
	Z = \frac{n}{n + \beta/\alpha + 1}.
	\]
	Cuando $n \to \infty$, $Z \to 1$ (confianza total en los datos propios).  Cuando $n \to 0$, $Z \to 0$ (confianza total en la experiencia colectiva).
\end{example}

\begin{center}
	\begin{tikzpicture}
		\begin{axis}[
			width=10cm, height=6cm,
			xlabel={$n$ (número de siniestros)},
			xlabel style={at={(axis description cs:0.5,-0.15)}},
			ylabel={$Z$ (factor de credibilidad)},
			ylabel style={at={(axis description cs:-0.15,0.5)}},
			domain=0:100,
			samples=100,
			axis lines=left,
			ymax=1.1,
			grid=both,
			grid style={gray!30},
			tick style={black},
		]
			\addplot[Accent, thick] {x/(x+10)};
			\addplot[AccentDark, dashed] coordinates {(0,1) (100,1)};
			\node[AccentDark] at (axis cs:80,0.85) {$Z \to 1$};
			\node[Accent] at (axis cs:10,0.3) {$K=10$};
		\end{axis}
	\end{tikzpicture}
\end{center}

\subsection{Conexión con la geometría}

La credibilidad bayesiana es una \textbf{proyección} en el espacio de distribuciones:

\begin{proposition}[Credibilidad como proyección]
	La prima de credibilidad bayesiana es la proyección $e$-exponencial de la distribución previa sobre la familia de distribuciones condicionales al historial del asegurado.
\end{proposition}

\section{Aprendizaje automático}

\subsection{Pérdida y geometría}

En aprendizaje automático, la función de pérdida codifica el ``costo'' de una predicción errónea.  Muchas pérdidas comunes son divergencias de Bregman.

\begin{theorem}[Pérdidas y Bregman]
	Sea $F$ una función convexa.  La pérdida $L(y, \hat{y}) = D_F(y \| \hat{y})$ es una pérdida de Bregman.  Casos especiales:
	\begin{itemize}
		\item $F(x) = x^2/2$: pérdida cuadrática (regresión).
		\item $F(x) = x\log x$: entropía cruzada (clasificación).
		\item $F(x) = e^x$: pérdida exponencial (boosting).
	\end{itemize}
\end{theorem}

\subsection{Regularización y geometría}

La regularización $L_2$ (ridge) es equivalente a una proyección $e$-exponencial sobre un conjunto de parámetros con norma acotada.

\begin{proposition}[Regularización como proyección]
	El problema de optimización regularizado:
	\[
	\min_\theta \frac{1}{n}\sum_{i=1}^n L(y_i, f_\theta(x_i)) + \lambda \|\theta\|^2
	\]
	es equivalente a la proyección $e$-exponencial de la verosimilitud sobre la bola $\|\theta\| \le r$ (donde $r$ depende de $\lambda$).
\end{proposition}

\section{Inferencia variacional}

\subsection{El problema de la aproximación bayesiana}

En inferencia bayesiana, el posterior $p(\theta | x)$ es a menudo intratable.  La \textbf{inferencia variacional} (VI) aproxima el posterior por una distribución más sencilla $q(\theta)$ minimizando la divergencia KL.

\begin{definition}[Inferencia variacional]
	El problema de VI es:
	\[
	\min_{q \in \mathcal{Q}} D_{\mathrm{KL}}(q \| p(\theta | x))
	\]
	donde $\mathcal{Q}$ es una familia paramétrica de distribuciones.
\end{definition}

\begin{theorem}[Descomposición ELBO]
	La divergencia KL se descompone como:
	\[
	D_{\mathrm{KL}}(q \| p(\theta | x)) = \underbrace{-\mathrm{ELBO}(q)}_{\text{evidencia inferior}} + \log p(x)
	\]
	donde la ELBO (Evidence Lower Bound) es:
	\[
	\mathrm{ELBO}(q) = \mathbb{E}_q[\log p(x, \theta)] - \mathbb{E}_q[\log q(\theta)].
	\]
	Minimizar la KL es equivalente a maximizar la ELBO.
\end{theorem}

\begin{proof}[Demostración]
	\begin{align*}
		D_{\mathrm{KL}}(q \| p(\theta|x))
		&= \mathbb{E}_q\!\left[\log \frac{q(\theta)}{p(\theta|x)}\right] \\
		&= \mathbb{E}_q[\log q(\theta)] - \mathbb{E}_q[\log p(\theta|x)] \\
		&= \mathbb{E}_q[\log q(\theta)] - \mathbb{E}_q[\log p(x,\theta)] + \log p(x) \\
		&= -\mathrm{ELBO}(q) + \log p(x).
	\end{align*}
\end{proof}

\begin{example}[VI para una Normal conjugada]
	Para $x | \theta \sim \mathcal{N}(\theta, 1)$ y $\theta \sim \mathcal{N}(0, \tau^2)$, el posterior es $\mathcal{N}(\mu_n, \sigma_n^2)$ con:
	\[
	\mu_n = \frac{n\bar{x}}{n + 1/\tau^2}, \quad \sigma_n^2 = \frac{1}{n + 1/\tau^2}.
	\]
	Si usamos VI con $q = \mathcal{N}(\mu, \sigma^2)$, la solución coincide con el posterior exacto (por conjugación).
\end{example}

\section{Métodos de gradiente natural}

\subsection{El problema de la optimización en espacio estadístico}

Los métodos de gradiente estándar (GD, SGD, Adam) ignoran la geometría del espacio de parámetros.  El \textbf{gradiente natural} corrige esto usando la métrica de Fisher.

\begin{definition}[Gradiente natural]
	El \textbf{gradiente natural} de una función $L(\theta)$ es:
	\[
	\tilde{\nabla} L(\theta) = G^{-1}(\theta) \nabla L(\theta)
	\]
	donde $G(\theta)$ es la matriz de información de Fisher.
\end{definition}

\begin{theorem}[Invariancia del gradiente natural]
	El gradiente natural es invariante bajo reparametrizaciones: si $\theta \to \phi(\theta)$, entonces $\tilde{\nabla} L$ se transforma como un vector contravariante.
\end{theorem}

\begin{proof}[Demostración (sketch)]
	Si $\phi = \phi(\theta)$, entonces $\nabla_\phi L = J^{-T} \nabla_\theta L$ y $G_\phi = J^T G_\theta J$, donde $J = \partial\phi/\\partial\theta$.  Entonces:
	\[
	\tilde{\nabla}_\phi L = G_\phi^{-1} \nabla_\phi L = (J^T G_\theta J)^{-1} J^T \nabla_\theta L = J^{-1} G_\theta^{-1} \nabla_\theta L = J^{-1} \tilde{\nabla}_\theta L.
	\]
	Esto es exactamente la transformación de un vector contravariante.
\end{proof}

\begin{example}[Gradiente natural para Bernoulli]
	Para $\mathrm{Bern}(p)$, la métrica de Fisher es $g = 1/(p(1-p))$.  Si la función de pérdida es $L(p) = -\log p(y|p)$, el gradiente natural es:
	\[
	\tilde{\nabla} L = p(1-p) \cdot \frac{\partial L}{\partial p}.
	\]
\end{example}

\begin{center}
	\begin{tikzpicture}
		\node[draw, rounded corners, fill=AccentLight!20, minimum width=2.5cm, minimum height=0.8cm] (gd) at (0,0) {Gradiente};
		\node[draw, rounded corners, fill=Accent!20, minimum width=2.5cm, minimum height=0.8cm] (ngd) at (4,0) {Gradiente Natural};
		\draw[->, thick, AccentDark] (gd) -- (ngd) node[midway, above] {$\times G^{-1}$};
		\node[below=0.3cm] at (gd.south) {Ignora geometría};
		\node[below=0.3cm] at (ngd.south) {Respeta geometría};
	\end{tikzpicture}
\end{center}

\section{Código Python: gradiente natural y VI}

\begin{lstlisting}[language=Python, style=PythonStyle]
import numpy as np

def natural_gradient_bernoulli(y_data, lr=0.01, n_steps=100):
    """Gradiente natural para estimar p en Bernoulli."""
    p = 0.5  # inicializacion
    history = [p]
    
    for _ in range(n_steps):
        # Gradiente de la log-verosimilitud
        grad = np.sum(y_data) / p - (len(y_data) - np.sum(y_data)) / (1 - p)
        
        # Metrica de Fisher
        fisher = 1 / (p * (1 - p))
        
        # Gradiente natural
        nat_grad = grad / fisher
        
        # Actualizacion
        p = p - lr * nat_grad
        p = np.clip(p, 0.01, 0.99)
        history.append(p)
    
    return p, history

# Verificacion
np.random.seed(42)
true_p = 0.7
data = np.random.binomial(1, true_p, size=200)
est_p, hist = natural_gradient_bernoulli(data, lr=0.1, n_steps=50)
print(f"Estimado: {est_p:.4f}, Verdadero: {true_p}")
print(f"Error: {abs(est_p - true_p):.4f}")
\end{lstlisting}

\begin{lstlisting}[language=Python, style=PythonStyle]
# Inferencia variacional para Normal conjugada
def variational_inference_normal(x_data, prior_mu=0, prior_var=1):
    """VI para modelo Normal-Normal."""
    n = len(x_data)
    x_bar = np.mean(x_data)
    
    # Posterior exacto (para comparacion)
    posterior_var = 1 / (n + 1/prior_var)
    posterior_mu = posterior_var * (n * x_bar + prior_mu/prior_var)
    
    # VI con familia Gaussiana
    # En este caso, VI coincide con el posterior exacto (conjugacion)
    q_mu = posterior_mu
    q_var = posterior_var
    
    # ELBO
    elbo = (-0.5 * np.log(2 * np.pi * q_var)
            - 0.5 * (q_var + (q_mu - x_bar)**2) * n
            + 0.5 * np.log(2 * np.pi * prior_var)
            + 0.5 * (q_mu**2 + q_var) / prior_var)
    
    return q_mu, q_var, elbo

# Verificacion
np.random.seed(42)
true_mu = 2.0
data = np.random.normal(true_mu, 1, size=50)
mu_q, var_q, elbo = variational_inference_normal(data, prior_mu=0, prior_var=1)
print(f"q_mu = {mu_q:.4f}, q_var = {var_q:.4f}")
print(f"ELBO = {elbo:.4f}")
\end{lstlisting}

\section{Ejercicios}

\subsection*{Credibilidad actuarial}
\begin{enumerate}
	\item Para un modelo Bühlmann con EPV $= 2$ y VHM $= 8$, calcule el factor de credibilidad $Z$ para $n = 20$.
	\item En el modelo Poisson--Gamma, muestre que $Z \to 1$ cuando $n \to \infty$.
	\item Para siniestros $\sim \mathrm{Poisson}(\lambda)$ con previa $\lambda \sim \mathrm{Gamma}(3, 1)$, calcule la prima de credibilidad para observaciones $3, 1, 4, 1, 5$.
	\item ¿Cuándo la credibilidad bayesiana coincide con la credibilidad Bühlmann?
\end{enumerate}

\subsection*{Aprendizaje automático}
\begin{enumerate}[resume]
	\item Muestre que la pérdida de entropía cruzada binaria es una divergencia de Bregman.
	\item Para regularización $L_2$, demuestre que la solución es la proyección $e$-exponencial sobre la bola de radio $r$.
	\item Calcule el gradiente natural de la función de pérdida cuadrática para una regresión lineal.
	\item Para una red neuronal con activación softmax, muestre que la salida es la proyección $m$-exponencial del vector de logits sobre el simplex.
\end{enumerate}

\subsection*{Inferencia variacional}
\begin{enumerate}[resume]
	\item Derive la ELBO para un modelo de mezcla de Gaussianas.
	\item Para $q = \mathcal{N}(\mu, \sigma^2)$ y posterior exacto $\mathcal{N}(\mu_n, \sigma_n^2)$, calcule $D_{\mathrm{KL}}(q \| p)$.
	\item Muestre que maximizar la ELBO es equivalente a minimizar la KL.
	\item Para un modelo lineal bayesiano, derive la solución VI cuando la previa es Gaussiana.
\end{enumerate}

\subsection*{Gradiente natural}
\begin{enumerate}[resume]
	\item Calcule el gradiente natural para $\mathrm{Poisson}(\lambda)$ con función de pérdida de entropía cruzada.
	\item Demuestre que el gradiente natural es invariante bajo reparametrizaciones.
	\item Para una red neuronal con activación softmax, muestre que el gradiente natural es equivalente al gradiente estándar multiplicado por la inversa de la matriz de Fisher.
	\item (Reto) Para una familia exponencial uniparamétrica, muestre que el paso de gradiente natural es $\Delta\eta = I(\eta)^{-1} \partial L / \partial\eta$, que es independiente de la escala del parámetro.
\end{enumerate}
